% ==========================================================
% Term Project Report
% ==========================================================

\documentclass{article}
\usepackage{spconf,amsmath,graphicx,hyperref}

% -------------------------
% Additional useful packages
% -------------------------
\usepackage{booktabs}
\usepackage{siunitx}
\usepackage{enumitem}
\usepackage{stfloats}
\usepackage{tikz}
\usetikzlibrary{arrows.meta, positioning}

% -------------------------
% Custom definitions
% -------------------------
\def\x{{\mathbf x}}
\def\L{{\cal L}}

% -------------------------
% Title and Author Info
% -------------------------
\title{Influence of ENSO in Indian and Sri Lankan Monsoon Rainfall}

\name{Abilaash Vijeyakumaran (Roll Number: 23035010628)} % TODO: fill in your name and roll number

\address{
Term Project Report - III\\
Trimester 9\\[0.5em]
B.Sc. (Honours) Data Science and Artificial Intelligence\\
Indian Institute of Technology Guwahati, India\\
\texttt{v.abilaash@op.iitg.ac.in} % TODO: fill in your email
}


\begin{document}
\ninept
\maketitle

% ==========================================================
% Abstract
% ==========================================================
\begin{abstract}
This study looks at how the El Ni\~{n}o-Southern Oscillation (ENSO) affected rainfall across
Sri Lanka and India from 1996 to 2025. It then applies the findings to the ongoing
El Ni\~{n}o event in 2026. Rather than being restricted to a single presumed monsoon season,
four zones such as, India Wet (Western Ghats), India Dry (Deccan), Sri Lanka Wet (South West),
and Sri Lanka Dry (North East), were analyzed independently using the Oceanic Ni\~{n}o Index
(ONI) and gridded CHIRPS precipitation. Correlation was examined at monthly resolution
throughout the entire annual cycle. This monthly-resolved analysis shows that both countries
have real, statistically significant ENSO sensitivity, but it is concentrated in discrete,
physically separate times. India's signal is strongest in August--October
(Dry zone August:$r=-0.562$, $p=0.001$); Sri Lanka shows a strong positive (wetter) signal
in October ($r=+0.594$, $p=0.0005$, the strongest single correlation in this study) followed
by a significant drying signal spanning February--April ($r=-0.37$ to $-0.51$ across zones and months).
The monthly decomposition reveals that the previous JJAS-only analysis's conclusion that
Sri Lanka had no ENSO sensitivity was an artifact of testing the incorrect season rather
than a lack of a true link. Applying these relationships to the current event -- an
El Ni\~{n}o Advisory as of August 2026, with Ni\~{n}o~3.4 forecast to reach $+3.4^{\circ}$C
by October and a 97\% chance of persistence through early spring 2027 -- projects
a substantially below-normal remainder of India's 2026 monsoon (most acute in August, 
$-3.5$ to $-4.2$~mm/day), an anomalously wet October 2026 for Sri Lanka ($+3$ to $+8$~mm/day,
conservatively capped at the historical maximum observed ONI to avoid extrapolation
beyond the historical record), and a multi-month drought risk for Sri Lanka spanning
February--April 2027 ($-1.5$ to $-4.1$~mm/day). These projections are broadly consistent
with, and provide historical statistical grounding for, public advisories issued by both
governments regarding this event.
\end{abstract}

% ==========================================================
% Keywords
% ==========================================================
\begin{keywords}
El Ni\~{n}o, ENSO, Oceanic Ni\~{n}o Index, Indian Monsoon, CHIRPS, Precipitation Variability
\end{keywords}

% ==========================================================
\section{Introduction}
% ==========================================================

The El Ni\~{n}o-Southern Oscillation (ENSO) is the dominant mode of interannual climate
variability on Earth, arising from coupled ocean-atmosphere interactions in the
equatorial Pacific. Its warm phase, El Ni\~{n}o, and cool phase, La Ni\~{n}a, are known
to influence weather patterns far beyond the Pacific basin through atmospheric
teleconnections. Among these, the relationship between ENSO and the South Asian summer
monsoon is one of the most extensively studied, given the monsoon's central importance
to agriculture, water resources, and the economies of India and Sri Lanka.

Understanding how strongly, and under what conditions, ENSO modulates monsoon rainfall
is important both for seasonal forecasting and for anticipating drought or flood risk
in the region. While the broad relationship (El Ni\~{n}o tends to suppress monsoon
rainfall) is well established in the climate literature, the strength of this
relationship can vary considerably by sub-region, depending on local topography and
rainfall regime.

This project quantifies the ENSO--monsoon relationship over India and Sri Lanka for the
period 1996--2025, using the Oceanic Ni\~{n}o Index (ONI) as the ENSO indicator and
gridded CHIRPS precipitation data as the rainfall record. The region is split into a
wet zone (Western Ghats) and a dry, rain-shadow zone (Deccan Plateau) to test whether
the strength of the relationship differs between areas with strong orographic rainfall
and areas more directly dependent on large-scale monsoon circulation.

Section~2 describes the problem formulation and objectives. Section~3 details the
methodology, including data sources and event-detection approach. Experimental results
are presented in Section~4, followed by discussion and conclusions in Sections~5
and~6.

% ==========================================================
\section{Problem Statement and Objectives}
% ==========================================================

\subsection{Problem Statement}
El Ni\~{n}o events are known to influence global weather patterns, but the strength and
consistency of this influence over the Indian/Sri Lankan monsoon region, and whether
it varies between the two countries and between wetter and drier sub-regions within
each, requires direct quantitative verification against observational data rather than
being assumed from the general literature alone. This project addresses two related
questions: \emph{does the Oceanic Ni\~{n}o Index show a statistically significant
relationship with monsoon-season rainfall over India and Sri Lanka between 1996 and
2025, and does the strength of this relationship differ (i) between India and Sri
Lanka, and (ii) between a wet and dry zone within each country?}

\subsection{Objectives}
\begin{enumerate}[noitemsep]
    \item To identify and classify El Ni\~{n}o, La Ni\~{n}a, and Neutral periods over
    1996--2025 using the Oceanic Ni\~{n}o Index (ONI).
    \item To obtain and process gridded precipitation data (CHIRPS) over India and Sri
    Lanka, and to define four representative sub-regions: a wet and dry zone within
    each country, kept separate rather than merged, so that national-level differences
    can be examined alongside zonal ones.
    \item To quantify the correlation between ONI and monsoon-season (JJAS) rainfall
    anomalies in each of the four zones, including statistical significance testing.
    \item To compare mean rainfall anomalies across ENSO phases (El Ni\~{n}o, La
    Ni\~{n}a, Neutral) using analysis of variance (ANOVA), for each zone.
    \item To determine whether India or Sri Lanka shows stronger ENSO sensitivity, and
    to interpret the physical basis for any observed difference.
\end{enumerate}

% ==========================================================
\section{Methodology / Approach}
% ==========================================================

This analysis was adapted from a group project originally conducted for Madagascar,
which examined the relationship between ENSO, precipitation, and vegetation (NDVI) over
the Southwest Indian Ocean region. The same overall pipeline -- ONI-based event
detection, gridded precipitation retrieval, zonal clipping, and phase-stratified
correlation analysis -- was re-applied to the India/Sri Lanka monsoon region as an
independent extension, given the well-documented and physically distinct
ENSO--monsoon teleconnection in this region.

\subsection{Overall Workflow}
The workflow proceeds in four stages: (1) ONI data acquisition and El Ni\~{n}o event
classification, (2) definition of the India/Sri Lanka study region and its four
sub-zones, (3) retrieval and spatial clipping of CHIRPS precipitation data, and (4)
statistical analysis relating ONI to zonal rainfall anomalies.

\subsection{Technical Details}

\textbf{ENSO classification.} The Oceanic Ni\~{n}o Index (ONI)~\cite{noaa_oni}, a
3-month running mean of sea surface temperature anomalies in the Ni\~{n}o~3.4 region,
was obtained from NOAA's Climate Prediction Center for 1996--2025. Following NOAA's operational
definition, a period was classified as an El Ni\~{n}o (La Ni\~{n}a) event when ONI
remained at or above $+0.5^{\circ}$C (at or below $-0.5^{\circ}$C) for at least five
consecutive overlapping 3-month seasons; all other periods were classified as Neutral.

\textbf{Precipitation data.} Daily gridded precipitation at $0.25^{\circ}$ resolution
was obtained from the CHIRPS v2.0 dataset~\cite{funk2015chirps} (Climate Hazards Group,
UC Santa Barbara) for 1996--2025 and clipped to the India and Sri Lanka national
boundaries (Natural Earth administrative boundaries~\cite{naturalearth}).

\textbf{Zone definition.} Four sub-regions were defined to represent contrasting
rainfall regimes within each country, kept separate rather than merged across
countries: an \emph{India wet zone} covering the Western Ghats and Kerala coastal belt,
an \emph{India dry zone} covering the Deccan Plateau rain-shadow region, a \emph{Sri
Lanka wet zone} covering the southwest of the island, and a \emph{Sri Lanka dry zone}
covering the north, east, and southeast of the island. Each zone was clipped to its
respective country's national boundary independently, rather than to a combined
India--Sri Lanka boundary, so that country-level comparisons would not be confounded by
cross-border averaging.

\textbf{Statistical analysis.} Monthly rainfall anomalies were computed relative to
each month's 1996--2025 climatology for both zones. Pearson correlation coefficients
between ONI and rainfall anomaly were computed (i) across all months and (ii)
restricted to the monsoon season (June--September, JJAS), since the ENSO--monsoon
teleconnection is expected to be strongest during this period. One-way ANOVA was used
to test whether mean JJAS rainfall anomalies differed significantly across the three
ENSO phases (El Ni\~{n}o, La Ni\~{n}a, Neutral).

\subsection{Mathematical Formulation}

\textbf{ONI and rainfall anomaly.} ONI is the 3-month running mean of sea surface
temperature anomaly in the Ni\~{n}o~3.4 region relative to a 1991--2020 climatology,
classified as El Ni\~{n}o ($\geq+0.5^{\circ}$C) or La Ni\~{n}a ($\leq-0.5^{\circ}$C) for
at least 5 consecutive overlapping seasons, and Neutral otherwise. For each zone, the
rainfall anomaly $A_t = P_t - \bar{P}_{c(t)}$, where $P_t$ is zone-averaged monthly
precipitation and $\bar{P}_{c(t)}$ is the climatological mean for calendar month $c(t)$.

\textbf{Pearson correlation.} The linear ONI--rainfall relationship is quantified by
\begin{equation}
r = \frac{\sum_i (x_i - \bar{x})(y_i - \bar{y})}
{\sqrt{\sum_i (x_i - \bar{x})^2}\sqrt{\sum_i (y_i - \bar{y})^2}}
\end{equation}
with $x_i=\text{ONI}_i$, $y_i=A_i$; significance assessed via a two-sided $t$-test with
$n-2$ degrees of freedom.

\textbf{Linear regression.} To translate a given ONI value into a projected rainfall
anomaly (used for the current-event outlook in
Section~\ref{sec:current_outlook}), an ordinary least-squares regression is fit for
each zone and month/season of interest:
\begin{equation}
\hat{A} = \beta_1 \cdot \text{ONI} + \beta_0
\end{equation}
where the slope and intercept minimize the sum of squared residuals,
\begin{equation}
\beta_1 = \frac{\sum_i (x_i-\bar{x})(y_i-\bar{y})}{\sum_i (x_i-\bar{x})^2}, \qquad
\beta_0 = \bar{y} - \beta_1 \bar{x}
\end{equation}
again with $x_i=\text{ONI}_i$ and $y_i=A_i$ restricted to the relevant month(s).
Predictions made at an input ONI value outside the range of $x_i$ used to fit the
model constitute extrapolation rather than interpolation, and are flagged as such
wherever they occur in this report (Section~\ref{sec:current_outlook}).

\textbf{One-way ANOVA.} To test whether mean anomaly differs across phases $k \in
\{\text{El Ni\~{n}o}, \text{La Ni\~{n}a}, \text{Neutral}\}$:
\begin{equation}
F = \frac{\sum_k n_k(\bar{A}_k - \bar{A})^2 \,/\, (K-1)}
{\sum_k \sum_{i \in k} (A_i - \bar{A}_k)^2 \,/\, (N-K)}
\end{equation}
where $\bar{A}_k$, $n_k$ are the mean and count in phase $k$, $\bar{A}$ the overall
mean, $K=3$, and $N$ the total sample size. A large $F$ indicates phase means differ
more than expected from within-phase variability alone.

% ==========================================================
\section{Experiments and Results}
% ==========================================================

\subsection{Experimental Setup}
The analysis used ONI data from NOAA CPC (1996--2025, monthly/ seasonal resolution) and
CHIRPS v2.0 daily precipitation data ($0.25^{\circ}$ resolution, aggregated to monthly
means) over the same period. All processing was performed in Python using
\texttt{xarray}, \texttt{rioxarray}, \texttt{geopandas}, and \texttt{scipy.stats}.
Statistical significance was assessed at the conventional $p<0.05$ threshold, with
$p<0.01$ and $p<0.001$ noted where applicable.

\subsection{Results}

Figure~\ref{fig:zones_map} shows the India/Sri Lanka study region with all four zones
delineated, kept separate by country. Prior to this national decomposition, the
whole-region (combined) same-month correlation was weak and non-significant
($r=-0.039$, $p=0.4625$), motivating the split presented below.

\begin{figure}
    \centering
    \includegraphics[width=\linewidth]{figures/zones_map.png}
    \caption{The four zones used in this analysis: India Wet (Western Ghats) and India
    Dry (Deccan Plateau), each clipped to India's boundary; and Sri Lanka Wet (southwest)
    and Sri Lanka Dry (north/east), each clipped to Sri Lanka's boundary independently.}
    \label{fig:zones_map}
\end{figure}

\subsubsection{Monthly-Resolved ENSO Sensitivity}
\label{sec:monthly}

An initial analysis restricted to the June--September (JJAS) monsoon season -- the
convention typically used for the Indian monsoon -- found strong ENSO sensitivity in
both Indian zones but none in either Sri Lankan zone. Since Sri Lanka's own monsoon
cycle (Yala/Maha) does not align with India's, this JJAS-only test was re-run at
monthly resolution across the full annual cycle for all four zones, testing the
ONI--rainfall anomaly correlation independently for each of the 12 calendar months
(Figure~\ref{fig:monthly_heatmap}).

\begin{figure}
    \centering
    \includegraphics[width=\linewidth]{figures/monthly_sensitivity_heatmap.png}
    \caption{Pearson correlation between ONI and rainfall anomaly, by calendar month
    and zone (1996--2025). Significance: * $p<0.05$, ** $p<0.01$, *** $p<0.001$.}
    \label{fig:monthly_heatmap}
\end{figure}

This reveals that Sri Lanka does have real ENSO sensitivity -- the JJAS-only result
was an artifact of testing the wrong season, not evidence of no relationship.
Table~\ref{tab:monthly_highlights} summarizes the strongest windows found for each
country.

\begin{table}[h]
\centering
\caption{Strongest monthly ENSO--rainfall correlations by country/zone.}
\label{tab:monthly_highlights}
\begin{tabular}{llcc}
\toprule
Country & Window & Zone & $r$ ($p$) \\
\midrule
India & Aug & Dry & $-0.562$ ($p{=}0.001$) \\
India & Aug & Wet & $-0.479$ ($p{=}0.007$) \\
India & Jul--Oct & both & $-0.36$ to $-0.56$, all sig. \\
Sri Lanka & Oct & Wet & $+0.594$ ($p{=}0.0005$) \\
Sri Lanka & Feb & Wet & $-0.507$ ($p{=}0.004$) \\
Sri Lanka & Feb--Apr & both & $-0.37$ to $-0.51$, all sig. \\
\bottomrule
\end{tabular}
\end{table}

Notably, India's own signal is not uniform across JJAS: June and July are individually
weak and non-significant ($r\approx-0.27$, $p\approx0.14$ in both zones), with the
pooled JJAS result driven mainly by August--October. With 48 individual month--zone
tests conducted, 2--3 significant results would be expected by chance alone at
$p<0.05$; the strongest findings above ($p<0.001$ for India August and Sri Lanka
October) are far too strong to be attributable to this, while single-star results
($p\approx0.02$--$0.05$) warrant more caution.

Figure~\ref{fig:india_srilanka_timeseries} shows ONI alongside country-averaged
rainfall for India and Sri Lanka.

\begin{figure}
    \centering
    \includegraphics[width=\linewidth]{figures/india_srilanka_timeseries.png}
    \caption{ONI (6-month smoothed, top panel) alongside country-averaged rainfall for
    India and Sri Lanka separately (6-month smoothed, bottom panel), 1996--2025, with El
    Ni\~{n}o event periods shaded in red. Prior to this national decomposition, the
    same-month correlation using a single whole-region (India+Sri Lanka combined)
    average was weak and non-significant ($r=-0.039$, $p=0.4625$), motivating the
    split presented here.}
    \label{fig:india_srilanka_timeseries}
\end{figure}

Table~\ref{tab:correlations} presents the original JJAS-restricted correlation test --
the initial analysis before the monthly decomposition above was performed. Both Indian
zones show a strong, highly significant negative correlation; neither Sri Lankan zone
reaches significance under this India-centric season definition, which (as
Section~\ref{sec:monthly} shows) reflects the choice of season rather than an absence
of Sri Lankan ENSO sensitivity.

\begin{table}[h]
\centering
\caption{Pearson correlation between ONI and rainfall anomaly, by country, zone, and
season.}
\label{tab:correlations}
\begin{tabular}{llccc}
\toprule
Country & Zone & Period & $r$ & $p$ \\
\midrule
India & Wet (Ghats) & All months & $-0.208$ & $0.0001$ \\
India & Wet (Ghats) & JJAS only  & $-0.408$ & $<0.0001$ \\
India & Dry (Deccan) & All months & $-0.219$ & $<0.0001$ \\
India & Dry (Deccan) & JJAS only  & $-0.392$ & $<0.0001$ \\
Sri Lanka & Wet (SW) & All months & $+0.077$ & $0.1424$ \\
Sri Lanka & Wet (SW) & JJAS only  & $+0.092$ & $0.3166$ \\
Sri Lanka & Dry (N/E) & All months & $-0.008$ & $0.8735$ \\
Sri Lanka & Dry (N/E) & JJAS only  & $-0.104$ & $0.2593$ \\
\bottomrule
\end{tabular}
\end{table}

Table~\ref{tab:anova} presents the ANOVA results and mean JJAS rainfall anomalies by
ENSO phase for all four zones. Both Indian zones show a highly significant phase effect
($p<0.001$); neither Sri Lankan zone does ($p>0.6$ in both cases).

\begin{table}[h]
\centering
\caption{Mean JJAS rainfall anomaly (mm/day) by ENSO phase, and one-way ANOVA results,
by country and zone.}
\label{tab:anova}
\begin{tabular}{llccc}
\toprule
Country & Zone & El Ni\~{n}o & La Ni\~{n}a & Neutral \\
\midrule
India & Wet (Ghats) & $-1.21$ & $+0.72$ & $+0.10$ \\
India & Dry (Deccan) & $-1.23$ & $+0.88$ & $+0.03$ \\
Sri Lanka & Wet (SW) & $+0.52$ & $-0.23$ & $-0.08$ \\
Sri Lanka & Dry (N/E) & $-0.06$ & $+0.18$ & $-0.06$ \\
\bottomrule
\end{tabular}
\vspace{0.3em}

\begin{tabular}{llcc}
\toprule
Country & Zone & $F$ & $p$ \\
\midrule
India & Wet (Ghats) & $8.640$ & $0.0003$ \\
India & Dry (Deccan) & $8.484$ & $0.0004$ \\
Sri Lanka & Wet (SW) & $0.296$ & $0.7445$ \\
Sri Lanka & Dry (N/E) & $0.415$ & $0.6614$ \\
\bottomrule
\end{tabular}
\end{table}

An additional correlation analysis between rainfall and vegetation
(NDVI)~\cite{vermote2019ndvi} anomalies, stratified by zone and ENSO phase (1-month
lag), is presented in Table~\ref{tab:ndvi}. In
this analysis, both Sri Lankan zones show moderate negative correlations under La
Ni\~{n}a (SW: $r=-0.26$; N/E: $r=-0.24$), while the India Wet zone shows the strongest
positive correlations of any zone under El Ni\~{n}o and La Ni\~{n}a alike ($r=+0.14$ and
$r=+0.16$ respectively); notably, none of these NDVI correlations reach the
$p<0.05$ threshold used elsewhere in this analysis, so they are reported as suggestive
patterns rather than confirmed relationships.

\begin{table}[h]
\centering
\caption{Rainfall--NDVI anomaly correlation (1-month lag), by country, zone, and ENSO
phase.}
\label{tab:ndvi}
\begin{tabular}{llcc}
\toprule
Zone & Phase & $r$ & $p$ \\
\midrule
India Dry (Deccan) & El Ni\~{n}o & $-0.10$ & -- \\
India Dry (Deccan) & La Ni\~{n}a & $+0.11$ & -- \\
India Dry (Deccan) & Neutral      & $+0.02$ & -- \\
India Wet (Ghats)  & El Ni\~{n}o & $+0.14$ & -- \\
India Wet (Ghats)  & La Ni\~{n}a & $+0.16$ & -- \\
India Wet (Ghats)  & Neutral      & $-0.02$ & -- \\
Sri Lanka Dry (N/E) & El Ni\~{n}o & $+0.01$ & -- \\
Sri Lanka Dry (N/E) & La Ni\~{n}a & $-0.24$ & -- \\
Sri Lanka Dry (N/E) & Neutral      & $-0.10$ & -- \\
Sri Lanka Wet (SW)  & El Ni\~{n}o & $-0.31$ & -- \\
Sri Lanka Wet (SW)  & La Ni\~{n}a & $-0.26$ & -- \\
Sri Lanka Wet (SW)  & Neutral      & $-0.12$ & -- \\
\bottomrule
\end{tabular}
\end{table}

% ==========================================================
\section{Discussion}
% ==========================================================

The central finding of this analysis is a clear asymmetry between India and Sri Lanka:
both Indian zones show a strong, highly significant, and physically consistent
ENSO--rainfall relationship, while neither Sri Lankan zone shows a statistically
significant relationship at all. This was not the expected outcome going in --
the two countries share the same broad monsoon system and are geographically close --
and it emerged only once the analysis moved from a single combined India+Sri Lanka
region to four independently-clipped, country-specific zones. Notably, the same-month
correlation using the whole-region average, prior to any decomposition, was weak and
non-significant ($r=-0.039$, $p=0.4625$); it
is only once the region is split by country and zone, and restricted to the monsoon
season, that the strong Indian signal becomes visible -- underscoring how easily a
real, strong relationship in one part of a region can be masked by averaging it
together with a part of the region where no such relationship exists.

\textbf{India.} Both Indian zones show a strong, highly significant negative
correlation between ONI and JJAS rainfall (Western Ghats: $r=-0.408$, $p<0.0001$;
Deccan: $r=-0.392$, $p<0.0001$), and both show a highly significant ANOVA across ENSO
phases ($F=8.640$, $p=0.0003$ and $F=8.484$, $p=0.0004$ respectively). The swing
between El Ni\~{n}o and La Ni\~{n}a years is substantial and nearly identical in both
zones ($\sim$1.9--2.1~mm/day), consistent with the well-documented finding that El
Ni\~{n}o suppresses the Indian summer monsoon. Notably, the wet and dry Indian zones
show very similar correlation strength once Sri Lanka is removed from the averaging --
suggesting the "dry zone is more sensitive" pattern seen in an earlier, combined-country
version of this analysis was partly an artifact of Sri Lanka diluting the wet-zone
signal, rather than a genuine India-internal topographic effect.

\textbf{Sri Lanka.} Neither Sri Lankan zone shows a statistically significant
relationship with ONI in any test ($p>0.6$ for both ANOVAs). The Sri Lanka wet zone's
correlation is even \emph{positive} ($r=+0.092$), opposite in sign to every Indian
zone, though not statistically distinguishable from zero. A plausible explanation is
that Sri Lanka, as a small island, is more strongly influenced by local-scale processes
-- land-sea breeze, its own central highlands, and dual exposure to both the southwest
and northeast monsoons -- than by the large-scale circulation shift that governs
India's response to ENSO; its smaller zones and fewer CHIRPS grid cells may also make
the signal noisier.

The seasonal trend analysis (JJA and DJF, all four zones) shows a modest positive
JJA trend across all four zones (India Wet: $+0.019$; India Dry: $+0.035$; Sri Lanka
Wet: $+0.051$; Sri Lanka Dry: $+0.021$~mm/day/yr) -- notably strongest in Sri Lanka's
wet zone despite its lack of ENSO sensitivity, suggesting this trend is driven by
factors other than ENSO. The DJF trend is flat to slightly negative in all zones.

The rainfall--NDVI correlation results (Table~\ref{tab:ndvi}) show a suggestive but
statistically unconfirmed pattern -- both Sri Lankan zones show moderate negative
correlations under La Ni\~{n}a, while India Wet shows moderate positive correlations
under both El Ni\~{n}o and La Ni\~{n}a -- but none reach significance, so these should
be treated as directions for future work rather than confirmed findings.

\textbf{Current Event Outlook (2026--27).}
\label{sec:current_outlook}
Since this historical analysis was
completed, ENSO conditions have shifted dramatically: NOAA CPC's ENSO Diagnostic
Discussion (3 August 2026)~\cite{cpc_ensodisc_aug2026} confirms an El Ni\~{n}o
Advisory, with the weekly Ni\~{n}o~3.4 index at $+2.3^{\circ}$C, forecast to reach
$+2.7^{\circ}$C (Aug), $+3.1^{\circ}$C (Sep), and $+3.4^{\circ}$C (Oct), and a 97\%
chance of persistence through early spring 2027. Applying the month-specific
regressions from Section~\ref{sec:monthly} to this trajectory:

\emph{India, remainder of 2026 monsoon.} Using NOAA's forecast ONI values directly
(within or near the regressions' observed range): August $-3.5$ to $-4.2$~mm/day,
September $-2.4$ to $-2.5$~mm/day, October $-2.1$~mm/day (both zones) -- a sharp
deficit early, gradually easing.

\emph{Sri Lanka, October 2026.} The forecast ONI ($+3.4$) exceeds any value observed
in the 1996--2025 record for this zone's regression (max observed: $2.28$); the naive
linear extrapolation ($+11.8$~mm/day, wet zone) is not considered credible. Capping the
input at the historically observed maximum gives a more defensible estimate of
$+7.9$~mm/day (wet zone) and $+3.0$~mm/day (dry zone) -- still a substantial positive
anomaly, consistent in direction and statistical strength ($r=+0.594$, $p=0.0005$)
with the strongest correlation found anywhere in this study, but with the true
magnitude genuinely uncertain given the event's unprecedented forecast strength.

\emph{Sri Lanka, February--April 2027.} Using the 1997--98 and 2015--16 events' actual
ONI trajectory at equivalent lead time as a historical analog (peak $\sim$2.4, well
within the regressions' observed range, so no capping needed): February $-2.6$ to
$-2.7$~mm/day, March $-2.4$ to $-4.1$~mm/day, April $-1.5$ to $-2.8$~mm/day across
zones -- a sustained, multi-month drought signal.

\textbf{Relevance to public advisories.} Sri Lankan government communications
regarding this event have suggested elevated rainfall around November 2026 followed by
drought conditions by February 2027. This analysis is broadly consistent with that
outlook, though the statistical signal is strongest and most confident one month
earlier than communicated (October rather than November -- November alone did not
reach significance in this dataset, $p=0.14$, likely due to limited single-month
sample size rather than absence of a real effect) and the drought signal appears to
extend beyond February through at least April. For India, this analysis provides
direct statistical support for expecting a below-normal remainder of the 2026 monsoon,
most acute in August.

\textbf{Limitations.} The 30-year record (8--9 El Ni\~{n}o events, $n=30$ per
individual calendar month) limits statistical power, particularly for single-month
tests. With 48 month--zone combinations tested, some weaker single-star results may
reflect chance rather than a genuine effect. The October Sri Lanka projection required
extrapolation beyond the observed ONI range, introducing additional uncertainty. Zone
boundaries used approximate bounding boxes rather than precise climatological
zonation, and other known drivers such as the Indian Ocean Dipole were not included.

% ==========================================================
\section{Conclusion and Future Work}
% ==========================================================

This project demonstrates that the Oceanic Ni\~{n}o Index is a statistically
significant and physically consistent predictor of rainfall variability over both
India and Sri Lanka between 1996 and 2025 -- a conclusion that only emerged after
moving beyond a single assumed monsoon season to a full monthly-resolved analysis.
An initial JJAS-only test found strong Indian sensitivity (Western Ghats: $r=-0.408$;
Deccan: $r=-0.392$, both $p<0.0001$) but no significant Sri Lankan relationship,
suggesting an asymmetry between the two countries. Decomposing by calendar month
revealed this was an artifact of testing the wrong season for Sri Lanka, not an
absence of sensitivity: Sri Lanka shows a strong positive (wetter) signal in October
($r=+0.594$, $p=0.0005$, the strongest single correlation in this entire study) and a
significant, multi-month drying signal spanning February--April ($r=-0.37$ to
$-0.51$). India's own signal, meanwhile, was found to be concentrated in
August--October rather than uniform across JJAS, with June and July individually
non-significant.

Applying these relationships to the actively-developing 2026 El Ni\~{n}o event
(confirmed El Ni\~{n}o Advisory as of August 2026, forecast to reach $+3.4^{\circ}$C
by October) projects a substantially below-normal remainder of India's 2026 monsoon,
an anomalously wet October 2026 for Sri Lanka, and a multi-month drought risk for Sri
Lanka spanning February--April 2027 -- projections broadly consistent with, and
providing historical statistical grounding for, public advisories issued by both
governments. This progression -- from a combined-region analysis showing no signal,
to a country-split JJAS analysis showing an apparent India/Sri Lanka asymmetry, to a
monthly-resolved analysis showing both countries are ENSO-sensitive in their own
distinct seasons -- illustrates how methodological choices about temporal and spatial
aggregation can determine whether a real underlying relationship is detected or
missed entirely.

\textbf{Future Scope.} This work can be extended by: (i) incorporating the Indian
Ocean Dipole as an additional predictor alongside ONI, particularly relevant given
Sri Lanka's dual Pacific/Indian Ocean exposure; (ii) building a predictive model
(e.g., logistic regression or ARIMA) to forecast above/below-normal rainfall by month,
moving from description toward genuine forecasting; (iii) refining zone boundaries
using climatological rainfall zonation rather than
bounding boxes; and (iv) revisiting the October Sri Lanka projection once the
current event's actual peak intensity is observed, to assess how well the
capped-extrapolation approach performed against ground truth.

% ==========================================================
\section{Artifacts and Demonstrations}
% ==========================================================

\begin{itemize}[noitemsep]
    \item \textbf{Code Repository:} \url{https://github.com/Abilaash17/enso-india-srilanka-monsoon} % TODO: add your repo link
    \item \textbf{Video Demo:} \url{https://youtu.be/ImG48W8fgYc} % TODO: update as needed
\end{itemize}

\textbf{Note:}
During development, code debugging, documentation editing (grammar refinement and technical clarity), were supported by AI-assisted tools such as Claude (Anthropic). The author's personal effort and knowledge are reflected in all design decisions, implementation choices, data source selections, and written material.


% ==========================================================
% References
% ==========================================================
\vfill\pagebreak
{\fontsize{8}{9}\selectfont
\bibliographystyle{IEEEbib}
\bibliography{refs}
}

\end{document}
